
\documentclass[11pt,a4paper]{article}

% -------------------------------------------------
% Encoding and typography
% -------------------------------------------------
\usepackage[T1]{fontenc}
\usepackage[utf8]{inputenc}
\usepackage{microtype}
\usepackage{amsmath}
\usepackage{amssymb}

% -------------------------------------------------
% Page layout and tables
% -------------------------------------------------
\usepackage[
    top=2.5cm,
    bottom=2.5cm,
    left=2.7cm,
    right=2.7cm
]{geometry}
\usepackage{booktabs}
\usepackage{array}
\usepackage{xcolor}

% -------------------------------------------------
% Links
% -------------------------------------------------
\usepackage[hidelinks]{hyperref}

% -------------------------------------------------
% Paragraph formatting
% -------------------------------------------------
\setlength{\parindent}{0pt}
\setlength{\parskip}{0.55em}

\begin{document}

% =================================================
% Header
% =================================================
\begin{flushleft}

    {\LARGE\textbf{Geospatial Foundation Models Benchmarking}}\\[0.3em]

    {\large\color{black!65}
    Research Reflections and Open Methodological Questions
    }\\[1.1em]

    \textbf{Stefano Infusini}
    \hfill
    University of Genoa -- DIBRIS\\[-0.1em]

    {\small\color{black!65}
    MSc Computer Science -- Artificial Intelligence
    }
    \hfill
    {\small September 2026}

\end{flushleft}

\vspace{-0.6em}
\rule{\textwidth}{0.4pt}
\vspace{1.0em}

\section*{Purpose}

This is a living document for short methodological observations that may become useful for the thesis.
The goal is not to produce a final methodology here, but to preserve the strongest ideas in a form
that can be reviewed quickly and progressively verified against primary sources.

A benchmark score should not be interpreted as a property of the architecture alone. In practice,

\[
\begin{aligned}
S=f(&\text{pretraining data, sensor, bands, resolution, task,}\\
    &\text{geography, adaptation, decoder, training regime}).
\end{aligned}
\]

Therefore, a leaderboard compares complete pretrained systems under particular conditions; it does
not by itself identify which component caused the performance difference.

% =================================================
\section{1. PANGAEA: GFMs do not dominate every downstream task}

PANGAEA shows that pretrained GFMs are not uniformly better than supervised task-specific baselines.
On HLS Burn Scars, using 100\% of the downstream labels, the reported mIoU values include:

\begin{center}
\begin{tabular}{lr}
\toprule
Model & HLS Burn Scars mIoU (\%) \\
\midrule
UNet baseline & 84.51 \\
Prithvi & 83.62 \\
CROMA & 82.42 \\
S12-MAE & 81.91 \\
S12-DINO & 81.72 \\
Scale-MAE & 76.68 \\
RemoteCLIP & 76.59 \\
\bottomrule
\end{tabular}
\end{center}

UNet is also among the top two models on six of the eleven datasets in PANGAEA Table 5
\cite{pangaea}. This is important because ``foundation model'' should not be interpreted as
``universally superior model.''

A more useful question is whether pretraining helps under specific conditions, especially when labels
are scarce or the downstream distribution differs from the training geography.

PANGAEA also reports 10\%, 50\% and 100\% label regimes. On HLS Burn Scars, for example,
SpectralGPT reaches 83.35 mIoU with 10\% of the labels, while UNet reaches 79.46
\cite{pangaea}. This suggests that the value of pretraining should be studied as a function of
label availability rather than only at full supervision.

\textbf{Thesis implication.}
A central comparison should be how performance degrades as the amount of geographically independent
supervision decreases.

% =================================================
\section{2. mIoU can hide the behaviour of the target class}

For class \(c\),

\[
\mathrm{IoU}_c=\frac{TP_c}{TP_c+FP_c+FN_c},
\]

and

\[
\mathrm{mIoU}=\frac{1}{C}\sum_{c=1}^{C}\mathrm{IoU}_c.
\]

For binary segmentation,

\[
\mathrm{mIoU}=
\frac{\mathrm{IoU}_{\mathrm{background}}+\mathrm{IoU}_{\mathrm{foreground}}}{2}.
\]

mIoU is much less sensitive to class imbalance than pixel accuracy because each class receives the
same weight. However, in applications such as burn-scar or flood mapping, the foreground class is
usually the scientifically important one. A high and stable background IoU can compress changes in
foreground performance.

A reported mIoU alone is therefore insufficient to reconstruct the foreground IoU. For binary EO
segmentation, the thesis should preferably report at least:

\[
\mathrm{mIoU},\qquad
\mathrm{IoU}_{\mathrm{foreground}},\qquad
\mathrm{IoU}_{\mathrm{background}},
\]

and, where useful, foreground precision, recall and F1/Dice.

% =================================================
\section{3. Benchmark fairness depends on what information each model can use}

Different GFMs expect different sensors and spectral bands. PANGAEA adapts downstream data to each
encoder: RGB-pretrained models may use only RGB channels from a multispectral dataset, while missing
expected bands can be zero-padded \cite{pangaea}. Supervised baselines can instead be trained from
scratch on all available downstream bands.

Therefore,

\[
\text{same benchmark dataset} \neq \text{same effective input information}.
\]

A performance difference may reflect both representation quality and access to different spectral
information.

For every model--dataset pair, the evaluation should record:

\begin{itemize}
    \item native or adapted input;
    \item sensor match with pretraining;
    \item spectral-band overlap;
    \item GSD/resolution similarity;
    \item temporal compatibility;
    \item required preprocessing or zero-padding.
\end{itemize}

A useful compatibility taxonomy is:

\begin{center}
\begin{tabular}{>{\raggedright\arraybackslash}p{3.6cm}
                >{\raggedright\arraybackslash}p{9.0cm}}
\toprule
Native-compatible & Only ordinary normalization/resizing is required. \\
Deterministically adaptable & A transparent fixed transformation is required, e.g. band selection. \\
Substantial adaptation & Missing channels/modalities or major transformations alter the model interface. \\
Not meaningful & Adaptation is so severe that the result no longer tests the released representation fairly. \\
\bottomrule
\end{tabular}
\end{center}

% =================================================
\section{4. Pretraining alignment is different from pretraining exposure}

\textbf{Pretraining--downstream alignment} means that pretraining and downstream data are similar in
sensor, bands, resolution, geography or acquisition conditions. This can legitimately improve transfer.
PANGAEA itself notes that Prithvi, pretrained on HLS imagery, performs strongly on HLS Burn Scars
\cite{pangaea}.

This is not automatically leakage.

\textbf{Pretraining exposure} is a stronger issue: the pretraining corpus may contain the same images,
locations, dates or near-duplicates as the downstream benchmark,

\[
D_{\mathrm{pretrain}}\cap D_{\mathrm{benchmark}}\neq\varnothing.
\]

Even without downstream labels, such exposure can complicate claims about geographic generalization.

For each selected model and benchmark, it would be useful to audit, where possible:

\begin{itemize}
    \item exact dataset overlap;
    \item geographic overlap;
    \item temporal overlap;
    \item sensor/spectral overlap;
    \item resolution similarity.
\end{itemize}

This should be treated as something to verify, not something to assume.

% =================================================
\section{5. So2Sat LCZ42: why it is interesting for classification}

So2Sat LCZ42 is a 17-class patch-classification benchmark for Local Climate Zones, with 10 built
classes and seven natural classes. The original benchmark contains 400,673 paired Sentinel-1/Sentinel-2
samples. Each sample is \(32\times32\) pixels at 10\,m, i.e. approximately \(320\times320\) m
\cite{so2sat}.

Its labels originate from manually delineated LCZ polygons. A Landsat-8 Random Forest is used during
annotation to produce a dense LCZ map and flag potentially inconsistent or missing training regions.
It is best interpreted as an annotation-debugging tool, not as independent ground truth. Final label
quality is further checked by human review and by a relabeling experiment involving ten remote-sensing
experts \cite{so2sat}.

The original split trains on 32 main cities plus ten add-on areas and evaluates on ten other cities.
Validation and test are geographically separated within those same held-out cities
\cite{so2sat}. This is useful for cross-city evaluation, but validation and test still share city-level
domains; a stricter protocol could use disjoint training, validation and test cities.

% =================================================
\section{6. So2Sat is difficult, but not close to random}

An overall accuracy of 50--60\% on So2Sat should not be interpreted like 50\% accuracy in balanced
binary classification.

For two balanced classes,

\[
P(\mathrm{correct})=\frac{1}{2}=50\%.
\]

For 17 equally likely classes,

\[
P(\mathrm{correct})=\frac{1}{17}\approx5.9\%.
\]

So historical So2Sat accuracies around 50--60\% are far above uniform random guessing. A majority-class
baseline should also be reported because the real class distribution is not perfectly balanced.

The dataset is difficult because several LCZ classes are semantically close, human experts themselves
show non-negligible disagreement, and the benchmark includes geographic transfer to unseen cities.

This difficulty is potentially useful. A benchmark saturated at 98--99\% has little power to distinguish
models or label-efficiency regimes. Conversely, a benchmark where all models collapse would suffer from
a floor effect. The useful regime is intermediate: difficult enough to preserve headroom, but not so
difficult that all models become indistinguishable.

Before fixing So2Sat as the final classification benchmark, modern GFMs should therefore be tested under
a consistent protocol to verify that it lies in this informative regime.

% =================================================
\section{7. So2Sat is not natively compatible with every multimodal GFM}

So2Sat supplies eight derived Sentinel-1 channels and ten Sentinel-2 bands \cite{so2sat}. The released
CROMA implementation expects two Sentinel-1 channels and twelve Sentinel-2 channels \cite{croma}.

Thus So2Sat may be a good classification benchmark while still being non-native for CROMA.

A deterministic adaptation could select the So2Sat SAR channels most closely corresponding to VV/VH
intensity/backscatter and explicitly handle missing optical bands. Such a result should be labelled as
an adapted setting rather than native CROMA input.

This motivates a broader rule: a benchmark should not force every model onto every dataset if the required
adaptation makes the comparison scientifically difficult to interpret.

% =================================================
\section{8. How to interpret cross-model results}

If

\[
FM_A > FM_B
\]

on one benchmark, the safe interpretation is:

\begin{quote}
Under this sensor, task, domain and adaptation configuration, \(FM_A\) transfers better than \(FM_B\).
\end{quote}

It does not automatically mean that \(FM_A\) has a better architecture or universally better
representations.

A useful experimental summary would combine performance with compatibility metadata:

\begin{center}
\begin{tabular}{lcccc}
\toprule
 & Classification & Segmentation & Regression & Change detection \\
\midrule
CROMA  & ? & ? & ? & ? \\
DOFA   & ? & ? & ? & ? \\
Prithvi& ? & ? & ? & ? \\
Clay   & ? & ? & ? & ? \\
\bottomrule
\end{tabular}
\end{center}

For each result, record whether input is native, sensor/bands match pretraining, GSD is similar,
geographic exposure is known, and adaptation is required.

This contextual information may be more informative than a single global leaderboard.

% =================================================
\section{9. Current thesis-level questions}

\begin{enumerate}
    \item Does GFM sample efficiency survive when spatial leakage is removed and training labels are
    geographically independent?

    \item How much of cross-model performance can be explained by pretraining--downstream alignment?

    \item Are model rankings stable when analysis is restricted to native-compatible model--dataset pairs?

    \item For binary segmentation, do foreground-only metrics change the conclusions obtained from mIoU?

    \item Does geographic OOD performance change when validation and test cities are also disjoint?

    \item Can pretraining exposure be audited well enough to distinguish transfer to unseen geographies
    from transfer to locations already encountered during pretraining?
\end{enumerate}

% =================================================
\section*{Working rule}

The objective should not be only to ask \emph{which GFM is best?}, but rather:

\emph{Under which data, supervision and geographic conditions does each pretrained system transfer well?}

% =================================================
\begin{thebibliography}{9}

\bibitem{pangaea}
V. Marsocci et al.,
\textit{PANGAEA: A Global and Inclusive Benchmark for Geospatial Foundation Models},
arXiv:2412.04204, 2024.
\href{https://arxiv.org/abs/2412.04204}{arXiv link}.

\bibitem{so2sat}
X. X. Zhu et al.,
\textit{So2Sat LCZ42: A Benchmark Data Set for the Classification of Global Local Climate Zones},
IEEE Geoscience and Remote Sensing Magazine, 8(3), 76--89, 2020.
DOI: 10.1109/MGRS.2020.2964708.

\bibitem{croma}
A. Fuller, K. Millard, and J. R. Green,
\textit{CROMA: Remote Sensing Representations with Contrastive Radar-Optical Masked Autoencoders},
NeurIPS, 2023.
\href{https://github.com/antofuller/CROMA}{Official repository}.

\end{thebibliography}

\end{document}
